\subsection{Proposed system}

\subsubsection{Description}

The proposed system is a comprehensive digital distribution platform and community hub for video game content, conceptually modeled after the industry-leading platform, Steam. Its predominant purpose is to serve as a complete, end-to-end ecosystem that connects content developers with consumers, facilitating not only the publishing and purchase of digital entertainment but also fostering a persistent, social community around those products.

\subsubsection{Proposed Application Overview \& Users}

The system is designed and configured around two primary actors with distinct roles and privileges, satisfying the Two-Sided Marketplace (B2B2C ecosystem) requirement, which concentrates on high-volume, user-friendly experiences, highlighting fast checkouts, personalized marketing, and robust data privacy.

The first actor is the \textbf{Player/Consumer}. These are the end-users who engage with the platform to discover new games and new genres, manage their personal collections, and bond socially, either through friendly conversation or constructive criticism. Their functions extend beyond simple purchasing to include building a persistent digital library, tracking in-game achievements, writing reviews, and fostering a network of friendships.

The second actor is the \textbf{Developer/Publisher}. These are the content creators or rights holders who utilize the platform as a sales and marketing channel. Their responsibilities consist of the initial listing of their games, delineating what the games entail in brief descriptions, managing post-launch content like downloadable content (DLC) or new seasonal rewards, updating game details, and ultimately tracking the sales performance of their titles.

The primary functions of the system can be grouped into five key areas. \textbf{User Management} handles the creation and maintenance of accounts for both Players and Developers, securing their access and storing their profile information. The \textbf{Storefront \& Purchasing} function is the commercial heart of the platform, allowing users to browse a compartmentalized catalog, manage a temporary shopping cart, and conduct financial transactions. \textbf{Library Management} is the persistent, personal space for each player, accurately tracking which specific games and DLCs they own, along with engagement metrics like total playtime. The \textbf{Social \& Community} function builds a layer of interaction on top of the catalog, enabling players to form friend connections, contribute reviews with ratings, comments, and constructive criticisms. In addition, the function allows players to track their progress in unlocking game-specific achievements, as well as view what achievements they have accomplished in specific games. Finally, \textbf{Publisher Functions} provide the backend tools for Developers to upload new products, manage game details and pricing, and create and associate DLC or seasonal updates like battle passes with their base games.

\subsubsection*{1.2.3 Entity Types, Attributes, and Relationships}

\noindent\textbf{1.2.3.1 Entity Types \& Attributes:}

The system's data model is centered on a hierarchy of user accounts, with \ent{USER} acting as a superclass. This entity captures the fundamental attributes common to all accounts, including a unique \attr{UserID} as its primary key, a chosen \attr{Username}, a verified \attr{Email} address, a secure \attr{PasswordHash}, and a \attr{JoinDate} marking their registration.

From this superclass, two distinct subclasses are specialized. The \ent{PLAYER} subclass inherits all \ent{USER} attributes and adds those specific to a consumer. The attributes of which comprise a \attr{WalletBalance}, which represents the user's stored funds and is a dynamic attribute updated with each transaction, and a \attr{ProfileBio}, a text field for personal expression on their profile page. The \ent{DEVELOPER} subclass, conversely, inherits the core \ent{USER} attributes but is augmented with business-oriented information. This includes a registered \attr{LegalName}---accommodating both individual creators and corporate entities. Furthermore, it incorporates a \attr{BankAccount} for receiving sales revenue, and importantly, a multivalued attribute, \attr{ContactPhone}, since a developer may maintain multiple contact numbers for miscellaneous reasons like customer support or business development

Perhaps, the most central product entity in Steam is, without a shadow of a doubt, \ent{GAME}. It is defined by a unique \attr{GameID} and several core descriptive attributes such as its \attr{Title}, a \attr{BasePrice}, and a \attr{ReleaseDate}. It also includes a derived attribute, \attr{EffectivePrice}, which is dynamically calculated at any given time by applying active discounts from associated promotions to the BasePrice. Its description is managed through a composite attribute, \attr{SystemRequirements}, which is itself composed of discrete elements like \attr{OS}, \attr{Processor}, \attr{Memory}, and \attr{Graphics}.

To properly facilitate game discovery, the system utilizes a classification entity: \ent{GENRE}. \ent{GENRE} is a strong entity representing the official, broad categories defined by the developer. It is uniquely identified by a \attr{GenreName} and includes a brief \attr{Description}. Because a single title can have the characteristics of multiple categories, a game is often linked to several genres; a survival horror title like \textit{Resident Evil Requiem} (2026) would be officially linked to both the "Action" and "Horror" genres, for instance.

Supporting the core game are several dependent entities. \ent{DLC} (Downloadable Content) is a weak entity, as its existence is entirely dependent on a specific base game and whether or not said base game performs well enough in sales to justify the development and release of a DLC. Its identity is partially determined by its relationship to a \attr{GameID} and is completed by its partial key, \attr{DLC\_Name}. It also has its own \attr{Price} and \attr{ReleaseDate}. Similarly, \ent{ACHIEVEMENT} is a weak entity dependent on a \ent{GAME}. It represents in-game challenges or milestones and is identified by its partial key, \attr{AchievementName}, its \attr{Description}, and also includes a descriptive \attr{Points} value.

\ent{REVIEW} is another weak entity that exists only in the context of a \ent{GAME}. It is identified by a partial key, \attr{ReviewID}, and contains the user's feedback, which includes a binary \attr{Rating} (representing a "Recommended" or "Not Recommended" thumbs-up/down), a detailed text \attr{Comment}, and the \attr{ReviewDate}.

Next, the \ent{TRANSACTION} entity is a strong entity representing a financial order. It is uniquely identified by a \attr{TransactionID} and records the details of a purchase, including the \attr{Date} when the game was initially purchased, the \attr{TotalAmount} paid for that specific game, and the \attr{PaymentMethod} utilized to make said transaction.

Another entity related to the purchasing workflow is the \ent{CART} entity, which acts as a temporary staging area for players before they finalize a purchase. It is uniquely identified by a \attr{CartID} and includes a \attr{LastUpdated} timestamp to track its recent activity and potentially recover abandoned carts. Unlike a permanent transaction record, a cart represents a player's current shopping intent, holding selected games in a temporary queue until checkout. At the moment of purchase, the system transfers these contents into a formal, permanent transaction record and subsequently empties the cart so it can be reused for future shopping sessions.

Last but not least, the \ent{PROMOTION} entity is a strong entity designed to handle dynamic pricing and storefront sales events, such as an annual summer sale or a specific publisher weekend. It is uniquely identified by a \attr{PromotionID} and contains core attributes, consisting of the event's \attr{Name}, the specific \attr{DiscountPercentage} being applied, as well as a \attr{StartDate} and \attr{EndDate} to automatically manage the lifecycle of the sale. Games are linked to these promotions over time, allowing the platform to offer temporary price reductions to consumers without permanently altering a game's standard base price.

\noindent\textbf{1.2.3.2 Relationships:}

The connections between these entities define the platform's logic. A \ent{DEVELOPER} \rel{Publishes} zero-to-many \ent{GAME}s, but each \ent{GAME} is published by exactly one \ent{DEVELOPER}, thus establishing a one-to-many (1:N) relationship.

The weak entities are tied to their owners through identifying relationships. A \ent{GAME} \rel{Has} many \ent{DLC} items, and because DLC is a weak entity as it cannot exist without the base game, this is an identifying 1:N relationship. The same pattern applies to \ent{ACHIEVEMENT}; a \ent{GAME} \rel{Has} \ent{ACHIEVEMENT}, forming an identifying 1:N relationship where the child cannot exist without the parent.

User actions are captured through several relationships. A \ent{PLAYER} \rel{Writes} many \ent{REVIEW}s, and because \ent{REVIEW} is a weak entity dependent on the player, this forms an identifying 1:N relationship. A \ent{PLAYER} also \rel{Makes} many \ent{TRANSACTION}s, with each transaction belonging to a single player, forming a regular 1:N relationship. Furthermore, a \ent{PLAYER} \rel{Has} a \ent{CART} in a strict one-to-one (1:1) relationship, where the \ent{CART} has total participation as it cannot exist without belonging to exactly one specific player.

The commercial and social cores are modeled with many-to-many (N:M) relationships. The purchase of games is captured by the \rel{Contains} relationship between \ent{TRANSACTION} and \ent{GAME}. A transaction can contain multiple games, and a game can appear in many transactions. This relationship requires its own attribute, \attr{PriceAtPurchase}, to accurately record the historical sale price independent of the game's current \attr{BasePrice}. Before that purchase is finalized, the \rel{Queues} N:M relationship links a \ent{CART} to a \ent{GAME}, representing the temporary intent to buy those specific titles. This relationship resolves into an associative entity that tracks which games are currently awaiting purchase in a given cart. Similarly, dynamic pricing is handled by the \rel{Applies\_To} N:M relationship between \ent{PROMOTION} and \ent{GAME}, meaning a game can participate in multiple sales events over time, and a specific sale features multiple games. This allows the system to calculate the \attr{EffectivePrice} of a game at any given time based on active promotions.

The player's library is modeled by the \rel{Owns\_Library} relationship between \ent{PLAYER} and \ent{GAME}. This N:M relationship tracks the games a player owns and is detailed with the \attr{PurchaseDate} and the dynamically updating \attr{TotalPlaytime}. Player achievements are tracked via the \rel{Unlocks} N:M relationship between \ent{PLAYER} and \ent{ACHIEVEMENT}, which includes the \attr{UnlockDate} attribute.

Finally, the social graph is represented by the recursive N:M relationship, \textbf{Friends\_With}, on the \ent{PLAYER} entity. This self-referential relationship allows a player to be friends with many other players and includes the \attr{DateAdded} to remember when the friendship was formed.